% ============================================================
%  Artículo Académico — Herramienta WHO_AM_I — MPU6050 / ESP32
%  Universidad del Cauca — Departamento de Ingeniería
%  Electrónica y Telecomunicaciones
%  Autor: Carlos Daniel Cuellar Antury
%  6 páginas — Formato doble columna, 10pt, Times New Roman
% ============================================================
\documentclass[10pt,twocolumn,a4paper]{article}

% --- Paquetes -----------------------------------------------------------
\usepackage[utf8]{inputenc}
\usepackage[T1]{fontenc}
\usepackage[spanish,es-tabla]{babel}
\usepackage{mathptmx}
\usepackage[top=2.5cm,bottom=2.5cm,left=2.5cm,right=2.5cm,
            columnsep=0.6cm]{geometry}
\usepackage{graphicx}
\usepackage{booktabs}
\usepackage{array}
\usepackage{listings}
\usepackage{xcolor}
\usepackage{hyperref}
\usepackage{caption}
\usepackage{float}
\usepackage{amsmath}
\usepackage{fancyhdr}
\usepackage{titlesec}
\usepackage{enumitem}
\usepackage{microtype}
\usepackage{url}

% --- Estilo de página ---------------------------------------------------
\pagestyle{fancy}
\fancyhf{}
\fancyhead[L]{\tiny Univ. del Cauca — Ing. Electrónica y Telecomunicaciones}
\fancyhead[R]{\tiny Cuellar Antury, C.D. (2024)}
\fancyfoot[C]{\small\thepage}
\renewcommand{\headrulewidth}{0.4pt}
\setlength{\headheight}{14pt}

% --- Formato de secciones -----------------------------------------------
\titleformat{\section}{\normalfont\bfseries\small}{\thesection.}{0.4em}{}
\titleformat{\subsection}{\normalfont\bfseries\footnotesize}{\thesubsection.}{0.4em}{}
\titleformat{\subsubsection}{\normalfont\itshape\footnotesize}{\thesubsubsection.}{0.4em}{}
\titlespacing{\section}{0pt}{5pt}{2pt}
\titlespacing{\subsection}{0pt}{4pt}{1pt}
\titlespacing{\subsubsection}{0pt}{3pt}{1pt}

% --- Listas compactas ---------------------------------------------------
\setlist{itemsep=1pt,topsep=2pt,parsep=0pt,partopsep=0pt}

% --- Colores para código ------------------------------------------------
\definecolor{codegray}{rgb}{0.96,0.96,0.96}
\definecolor{codeblue}{rgb}{0.0,0.0,0.55}
\definecolor{codegreen}{rgb}{0.0,0.45,0.0}
\definecolor{codepurple}{rgb}{0.45,0.0,0.45}
\definecolor{codered}{rgb}{0.6,0.0,0.0}

\lstdefinestyle{arduino}{
  language=C++,
  backgroundcolor=\color{codegray},
  basicstyle=\ttfamily\scriptsize,
  keywordstyle=\color{codeblue}\bfseries,
  commentstyle=\color{codegreen}\itshape,
  stringstyle=\color{codepurple},
  numberstyle=\tiny\color{gray},
  numbers=left,
  stepnumber=1,
  numbersep=4pt,
  breaklines=true,
  showspaces=false,
  showstringspaces=false,
  frame=single,
  rulecolor=\color{black!20},
  tabsize=2,
  captionpos=b,
  aboveskip=4pt,
  belowskip=2pt,
  morekeywords={uint8\_t,Wire,Serial}
}

\lstdefinestyle{output}{
  basicstyle=\ttfamily\scriptsize,
  backgroundcolor=\color{black!5},
  frame=single,
  numbers=none,
  captionpos=b,
  aboveskip=3pt,
  belowskip=2pt,
  breaklines=true
}

\hypersetup{colorlinks=true,linkcolor=blue,urlcolor=blue,citecolor=blue,
            breaklinks=true}
\captionsetup{font=scriptsize,labelfont=bf,skip=3pt}

% ========================================================================
\begin{document}

% -----------------------------------------------------------------------
%  PORTADA — TÍTULO, AUTOR, AFILIACIÓN, ABSTRACT (una columna)
% -----------------------------------------------------------------------
\twocolumn[{%
\begin{center}
  {\large\bfseries Herramienta de Diagnóstico de Sensor MPU6050 mediante\\[2pt]
  Identificación de ID Interno en Plataforma ESP32}\\[6pt]
  {\normalsize Carlos Daniel Cuellar Antury}\\[2pt]
  {\small Departamento de Ingeniería Electrónica y Telecomunicaciones,
  Universidad del Cauca, Popayán, Colombia}\\[2pt]
  {\footnotesize\url{https://github.com/Cdcuellar/Sistemas-embebidos---Actividades-de-la-vida-diaria}}\\[6pt]
  \rule{\linewidth}{0.5pt}\\[4pt]
\end{center}
\begin{small}
\noindent\textbf{Resumen} — Este artículo presenta el diseño, implementación y
validación de una herramienta de diagnóstico de hardware para el sensor MPU6050
sobre la plataforma ESP32 DevKit V1, desarrollada en el marco del proyecto de
sistemas wearables para monitoreo de Actividades de la Vida Diaria (AVD) en la
Universidad del Cauca. La herramienta lee el registro \texttt{WHO\_AM\_I}
(dirección \texttt{0x75}) mediante el protocolo I2C utilizando la biblioteca
nativa \texttt{Wire.h} del entorno Arduino IDE, con pines GPIO~21 (SDA) y
GPIO~22 (SCL), a una velocidad serial de 115\,200~baudios. El resultado
experimental arroja el valor \texttt{0x72} como identificador de silicio, que
difiere del \texttt{0x68} estándar indicado en la hoja de datos oficial de
InvenSense/TDK. Este hallazgo es crítico para la correcta configuración del
driver de bajo nivel en ESP-IDF. Se presentan tablas de especificaciones,
análisis del flujo I2C, listado del código fuente comentado y tabla de
diagnóstico de resultados posibles.

\smallskip
\noindent\textbf{Palabras clave} — MPU6050, WHO\_AM\_I, I2C, ESP32, diagnóstico
de hardware, sistemas embebidos, wearables, actividades de la vida diaria.

\smallskip
\noindent\textbf{Abstract} — This paper presents the design, implementation, and
validation of a hardware diagnostic tool for the MPU6050 sensor on the ESP32
DevKit V1 platform, developed within the wearable systems project for Activities
of Daily Living (ADL) monitoring at Universidad del Cauca. The tool reads the
\texttt{WHO\_AM\_I} register (address \texttt{0x75}) via I2C using the native
\texttt{Wire.h} library in the Arduino IDE environment, with pins GPIO~21 (SDA)
and GPIO~22 (SCL) at 115\,200~baud. The experimental result yields \texttt{0x72}
as the silicon identifier, differing from the standard \texttt{0x68} indicated
in the official InvenSense/TDK datasheet. This finding is critical for the
correct configuration of the low-level driver in ESP-IDF. Specification tables,
I2C flow analysis, annotated source code listing, and a diagnostic table for
possible outcomes are presented.

\smallskip
\noindent\textbf{Keywords} — MPU6050, WHO\_AM\_I, I2C, ESP32, hardware
diagnostics, embedded systems, wearables, activities of daily living.
\end{small}
\vspace{4pt}\rule{\linewidth}{0.5pt}\vspace{4pt}
}]

% -----------------------------------------------------------------------
%  1. INTRODUCCIÓN
% -----------------------------------------------------------------------
\section{Introducción}

El envejecimiento poblacional y el auge de la telemedicina han impulsado
el desarrollo de sistemas embebidos para el monitoreo continuo y
no intrusivo de Actividades de la Vida Diaria (AVD). Estos sistemas
—conocidos como wearables— requieren unidades de medición inercial
(IMU) capaces de capturar patrones de movimiento humano con alta
fidelidad, bajo consumo energético y reducido perfil físico~\cite{chen2012}.

El sensor MPU6050 de InvenSense/TDK se ha consolidado como una de las
opciones más extendidas en aplicaciones de investigación y prototipado
wearable gracias a su integración de acelerómetro y giroscopio triaxiales
en un único encapsulado de 4$\times$4~mm, su interfaz I2C estándar y su
bajo precio de mercado~\cite{invensense2013}. No obstante, la masificación
del mercado de componentes electrónicos ha generado una proliferación de
módulos MPU6050 que, si bien son compatibles en pinout, corresponden a
revisiones de silicio distintas a las especificadas en la documentación
oficial. Esta situación origina un problema práctico y crítico en el
desarrollo de firmware: la inicialización del sensor puede fallar
silenciosamente cuando el driver verifica la identidad del dispositivo
mediante el registro \texttt{WHO\_AM\_I} y obtiene un valor inesperado.

En el contexto del proyecto \textit{Sistema de Monitoreo Vestible para
Actividades de la Vida Diaria}, desarrollado en el Departamento de
Ingeniería Electrónica y Telecomunicaciones de la Universidad del
Cauca~\cite{repo_github}, el Nodo~2 (encargado del registro de
movimiento) utiliza el MPU6050 sobre la plataforma ESP32 DevKit~V1.
Durante la etapa de integración con el framework ESP-IDF, se detectó
que la librería de referencia rechazaba el sensor por discrepancia en
el valor de \texttt{WHO\_AM\_I}. Como respuesta, se diseñó e implementó
una herramienta de diagnóstico minimalista en Arduino IDE que permite
determinar el valor real del identificador de silicio antes de proceder
con el firmware de producción.

El presente artículo documenta el diseño, implementación y resultados
de esta herramienta, constituyendo un recurso técnico reproducible para
cualquier desarrollador que enfrente la misma problemática.

\subsection{Contexto del Proyecto Wearable}

El sistema wearable al que pertenece esta herramienta es una red de
nodos embebidos heterogéneos cuya función colectiva es capturar señales
biométricas y de movimiento para el reconocimiento automático de AVD.
El repositorio público~\cite{repo_github} contiene múltiples
subsistemas: lectura de frecuencia cardíaca (Nodo~1), detección de
movimiento con MPU6050 (Nodo~2), pantalla táctil ESP32-S3 (Nodo~3) y
detección de dirección MAC. Esta herramienta fue desarrollada
específicamente para resolver el problema de inicialización del Nodo~2.

\subsection{Organización del Artículo}

El resto del artículo se organiza como sigue: la Sección~\ref{sec:rw}
revisa trabajos relacionados; la Sección~\ref{sec:teoria} presenta el
marco teórico sobre el MPU6050, el protocolo I2C y el registro
\texttt{WHO\_AM\_I}; la Sección~\ref{sec:metodo} describe la metodología
y el diseño experimental; la Sección~\ref{sec:resultados} presenta los
resultados; la Sección~\ref{sec:discusion} ofrece la discusión;
finalmente, la Sección~\ref{sec:conclusiones} expone conclusiones y
trabajo futuro.

% -----------------------------------------------------------------------
%  2. TRABAJOS RELACIONADOS
% -----------------------------------------------------------------------
\section{Trabajos Relacionados}
\label{sec:rw}

El uso del MPU6050 en sistemas de monitoreo de AVD es ampliamente
documentado en la literatura. Anguita \textit{et al.}~\cite{anguita2012}
emplearon el MPU6050 en un teléfono inteligente para reconocimiento de
actividades físicas (caminar, subir escaleras, sentarse), logrando una
tasa de clasificación superior al 96\%. Su trabajo evidencia la
importancia de la calibración y verificación del sensor antes de
adquirir datos.

En el ámbito de sistemas embebidos de bajo nivel, la problemática de
módulos con identificadores de silicio no estándar ha sido reportada
en comunidades de desarrollo como Stack Overflow y los foros oficiales
de Espressif. Kolban~\cite{kolban2017} señala en su obra de referencia
sobre ESP32 que la verificación del registro de identidad del dispositivo
debe hacerse antes de asumir compatibilidad con cualquier driver I2C,
especialmente en hardware de terceros.

En sistemas IDF, la función de inicialización del driver MPU6050 en
proyectos de código abierto (p.ej. ESP-IDF MPU driver de hideakitai)
implementa una verificación de \texttt{WHO\_AM\_I} con máscara de bits
que solo acepta \texttt{0x68} o \texttt{0x70}, rechazando el valor
\texttt{0x72} presente en el hardware del presente proyecto. Este
comportamiento fue la motivación directa para el desarrollo de la
herramienta de diagnóstico aquí descrita.

A diferencia de los trabajos citados, la herramienta propuesta se centra
específicamente en el diagnóstico del valor de \texttt{WHO\_AM\_I} como
paso previo obligatorio a la integración del firmware, con un código
mínimo, sin dependencias externas y con documentación detallada de los
posibles resultados.

% -----------------------------------------------------------------------
%  3. MARCO TEÓRICO
% -----------------------------------------------------------------------
\section{Marco Teórico}
\label{sec:teoria}

\subsection{El Sensor MPU6050}

El MPU6050 es una unidad de medición inercial (IMU) de seis grados de
libertad (6-DOF) fabricada por InvenSense, actualmente división de
TDK Corporation. Integra en un único encapsulado QFN de 4$\times$4~mm:
(a) un acelerómetro triaxial MEMS de rango configurable, (b) un
giroscopio triaxial MEMS de rango configurable, y (c) un procesador
de movimiento digital (DMP) con 1~KB de FIFO~\cite{invensense2013}.

\begin{table}[H]
\centering
\caption{Especificaciones principales del MPU6050}
\label{tab:specs}
\scriptsize
\begin{tabular}{p{3.0cm}p{3.5cm}}
\toprule
\textbf{Parámetro} & \textbf{Valor} \\
\midrule
Interfaz de comunicación & I2C (hasta 400~kHz) \\
Tensión de alimentación & 2,375 -- 3,46~V \\
Rango acelerómetro & $\pm$2/$\pm$4/$\pm$8/$\pm$16~g \\
Rango giroscopio & $\pm$250 a $\pm$2000~°/s \\
Resolución ADC & 16 bits \\
Temperatura operación & $-40$ a $+85$~°C \\
Dirección I2C (AD0=0) & 0x68 \\
Dirección I2C (AD0=1) & 0x69 \\
Registro WHO\_AM\_I & 0x75 (esperado: 0x72*) \\
Consumo activo & 3,9~mA \\
Encapsulado & QFN 4$\times$4 mm \\
\bottomrule
\multicolumn{2}{l}{\tiny *Valor detectado en hardware del proyecto.} \\
\end{tabular}
\end{table}

La comunicación con el host se realiza mediante I2C estándar. El sensor
dispone de 8 rangos de sensibilidad para el acelerómetro y el
giroscopio, configurables mediante registros de 8 bits en el mapa de
memoria interno. Los datos de cada eje son registros de 16 bits en
complemento a dos, organizados en pares de registros consecutivos de
alta y baja resolución.

\subsection{El Protocolo I2C}

\subsubsection{Generalidades}

I2C (\textit{Inter-Integrated Circuit}) es un protocolo de comunicación
serie síncrono de dos hilos desarrollado por Philips Semiconductors en
1982 y estandarizado por NXP~\cite{nxp_i2c}. Opera con una línea de
datos bidireccional (SDA, \textit{Serial Data Line}) y una línea de
reloj unidireccional (SCL, \textit{Serial Clock Line}), ambas con
resistencias de pull-up hacia VCC. La topología es maestro-esclavo,
con soporte para múltiples dispositivos en el mismo bus.

\subsubsection{Modos de Operación}

El estándar I2C define cuatro modos de velocidad de transferencia:
Standard (100~kHz), Fast (400~kHz), Fast-Plus (1~MHz) y High-Speed
(3,4~MHz). En el presente proyecto se emplea el modo Standard a 100~kHz,
velocidad predeterminada de la biblioteca \texttt{Wire.h} para ESP32.

\subsubsection{Estructura de la Trama I2C}

Una transacción I2C de escritura+lectura (utilizada para acceder a
registros internos del MPU6050) comprende las siguientes fases:
\begin{enumerate}
  \item \textbf{START}: El maestro toma control del bus.
  \item \textbf{Dirección del esclavo} (7 bits) + bit W=0.
  \item \textbf{ACK} del esclavo.
  \item \textbf{Dirección del registro} a leer (8 bits).
  \item \textbf{ACK} del esclavo.
  \item \textbf{REPEATED START}: Reinicio sin liberar el bus.
  \item \textbf{Dirección del esclavo} (7 bits) + bit R=1.
  \item \textbf{ACK} del esclavo.
  \item \textbf{Dato leído} por el maestro (8 bits).
  \item \textbf{NACK} + \textbf{STOP}.
\end{enumerate}

\subsubsection{I2C en Arduino para ESP32}

En el entorno Arduino para ESP32, la biblioteca \texttt{Wire.h} abstrae
completamente el protocolo I2C~\cite{wire_arduino}. La inicialización
se realiza con \texttt{Wire.begin(SDA, SCL)}, especificando los números
de pin GPIO. Para el ESP32 DevKit V1, los pines por defecto son GPIO~21
para SDA y GPIO~22 para SCL.

\subsection{El Registro WHO\_AM\_I (0x75)}

El registro \texttt{WHO\_AM\_I} es un campo de solo lectura de 8 bits,
ubicado en la dirección de registro \texttt{0x75} del mapa de memoria
del MPU6050. Su función es identificar de manera unívoca el modelo del
dispositivo conectado en el bus I2C, lo que permite al firmware de host
verificar que el sensor es el esperado antes de inicializar los
registros de configuración.

Según la hoja de datos oficial~\cite{invensense2013}, el valor de este
registro para el MPU6050 es \texttt{0x68} (coincidente con la dirección
I2C de esclavos AD0=0, aunque son conceptualmente independientes). Sin
embargo, como resultado de la proliferación de clones y variantes de
hardware, en la práctica se han documentado los valores indicados en
la Tabla~\ref{tab:who}.

\begin{table}[H]
\centering
\caption{Valores del registro WHO\_AM\_I documentados}
\label{tab:who}
\scriptsize
\begin{tabular}{ccp{2.8cm}}
\toprule
\textbf{Valor HEX} & \textbf{Dispositivo} & \textbf{Observación} \\
\midrule
0x68 & MPU6050 original & Según datasheet oficial \\
0x72 & MPU6050 clon/rev. & Detectado en el proyecto \\
0x98 & MPU6000/variante & Rev. alternativa conocida \\
0x70 & MPU6500 & IMU de generación posterior \\
\bottomrule
\end{tabular}
\end{table}

La discrepancia entre el valor esperado y el valor real puede causar
que los drivers de bajo nivel en ESP-IDF rechacen el sensor durante la
inicialización, incluso cuando el hardware está físicamente funcional
y correctamente conectado.

\subsection{La Plataforma ESP32 DevKit V1}

El ESP32 (Espressif Systems) es un microcontrolador de doble núcleo
Xtensa LX6 a hasta 240~MHz, con 520~KB de SRAM, Wi-Fi 802.11 b/g/n,
Bluetooth 4.2, y soporte para múltiples interfaces de comunicación
(I2C, SPI, UART, I2S, CAN)~\cite{espressif2023}. La placa de desarrollo
\textit{DevKit V1} es la referencia estándar de la comunidad de Arduino
para ESP32, con 30 pines GPIO expuestos y regulador de 3.3~V integrado.

Los pines GPIO~21 (SDA) y GPIO~22 (SCL) son los asignados por defecto
para I2C en el núcleo Arduino ESP32, y son los utilizados en la
herramienta de diagnóstico aquí descrita. Ambos pines admiten pull-up
interno configurable por software, aunque en aplicaciones I2C es
recomendable usar resistencias externas de 4,7~k$\Omega$.

% -----------------------------------------------------------------------
%  4. METODOLOGÍA Y DISEÑO EXPERIMENTAL
% -----------------------------------------------------------------------
\section{Metodología y Diseño Experimental}
\label{sec:metodo}

\subsection{Configuración de Hardware}

El esquema de conexión entre el ESP32 DevKit V1 y el módulo MPU6050 se
muestra en la Figura~\ref{fig:conexion}. La alimentación del módulo se
toma directamente del pin de 3.3~V del ESP32. La Tabla~\ref{tab:pins}
detalla la asignación de pines y parámetros del sistema.

\begin{table}[H]
\centering
\caption{Configuración de pines y parámetros del sistema}
\label{tab:pins}
\scriptsize
\begin{tabular}{p{1.8cm}p{1.5cm}p{2.8cm}}
\toprule
\textbf{Señal} & \textbf{Pin ESP32} & \textbf{Descripción} \\
\midrule
SDA & GPIO 21 & Datos I2C bidireccional \\
SCL & GPIO 22 & Reloj I2C \\
VCC & 3.3V & Alimentación módulo \\
GND & GND & Tierra común \\
AD0 & GND & Dir. I2C = 0x68 \\
\midrule
\textbf{Parámetro} & \multicolumn{2}{l}{\textbf{Valor}} \\
\midrule
Velocidad I2C & \multicolumn{2}{l}{100 kHz (Standard)} \\
Velocidad serial & \multicolumn{2}{l}{115\,200 baudios} \\
Dirección I2C & \multicolumn{2}{l}{0x68 (AD0=GND)} \\
Registro leído & \multicolumn{2}{l}{0x75 (WHO\_AM\_I)} \\
\bottomrule
\end{tabular}
\end{table}

\begin{figure}[H]
\centering
\scriptsize
\begin{verbatim}
 +--------------------+    +------------------+
 |   ESP32 DevKit V1  |    |     MPU6050      |
 |                    |    |                  |
 |  GPIO21 (SDA) -----+----+-- SDA            |
 |  GPIO22 (SCL) -----+----+-- SCL            |
 |  3.3V         -----+----+-- VCC            |
 |  GND          -----+----+-- GND            |
 |                    |    +-- AD0 ---[GND]   |
 +--------------------+    +------------------+
   I2C Master (ESP32)        I2C Slave (0x68)
\end{verbatim}
\caption{Diagrama de conexión I2C: ESP32 DevKit V1 — MPU6050}
\label{fig:conexion}
\end{figure}

\subsection{Entorno de Desarrollo Software}

\begin{itemize}
  \item \textbf{IDE}: Arduino IDE (versión 2.x, con soporte ESP32)
  \item \textbf{Placa objetivo}: ESP32 Dev Module (esp32:esp32:esp32doit-devkit-v1)
  \item \textbf{Biblioteca I2C}: \texttt{Wire.h} nativa del núcleo ESP32;
        no se requieren librerías externas adicionales.
  \item \textbf{Monitor Serie}: 115\,200~baudios, finalización de línea
        \textit{Newline}.
  \item \textbf{Sistema Operativo}: Compatible con Windows, Linux,
        macOS.
  \item \textbf{Compilador}: GCC Xtensa (incluido en el paquete del
        núcleo ESP32 para Arduino).
\end{itemize}

\subsection{Descripción del Código Fuente}

El Listado~\ref{lst:code} muestra el código fuente completo de la
herramienta, disponible en el repositorio GitHub del
proyecto~\cite{repo_github}. El programa sigue una arquitectura
deliberadamente minimalista: toda la lógica de diagnóstico se ejecuta
en la función \texttt{setup()}, que corre una única vez al encendido
del microcontrolador. La función \texttt{loop()} permanece vacía para
evitar lecturas repetitivas innecesarias.

\begin{lstlisting}[style=arduino,
  caption={herramienta\_who\_i\_am\_id\_interno\_arduino\_ide.ino},
  label=lst:code]
#include <Wire.h>

void setup() {
  // Iniciar I2C: SDA=GPIO21, SCL=GPIO22
  Wire.begin(21, 22);
  // Iniciar comunicacion serial a 115200 baudios
  Serial.begin(115200);

  Serial.println("Leyendo WHO_AM_I...");

  // Iniciar transmision al esclavo 0x68
  Wire.beginTransmission(0x68);
  // Escribir direccion del registro WHO_AM_I
  Wire.write(0x75);
  // Repeated START (no liberar el bus)
  Wire.endTransmission(false);

  // Solicitar 1 byte de lectura al esclavo
  Wire.requestFrom(0x68, 1);

  if (Wire.available()) {
    uint8_t val = Wire.read();
    Serial.print("WHO_AM_I: 0x");
    Serial.println(val, HEX);
  } else {
    Serial.println("No se pudo leer el sensor");
  }
}

// Sin operaciones en el bucle principal
void loop() {}
\end{lstlisting}

\noindent{\tiny Código fuente en GitHub:\\
\url{https://github.com/Cdcuellar/Sistemas-embebidos---Actividades-de-la-vida-diaria/blob/main/herramienta_who_i_am_id_interno_arduino_ide/herramienta_who_i_am_id_interno_arduino_ide.ino}}

\medskip
\noindent{\tiny Notas de depuración (README.txt):\\
\url{https://github.com/Cdcuellar/Sistemas-embebidos---Actividades-de-la-vida-diaria/blob/main/herramienta_who_i_am_id_interno_arduino_ide/README.txt}}

\subsection{Análisis Línea a Línea del Código}

\subsubsection{Inicialización del Bus I2C}
La llamada \texttt{Wire.begin(21, 22)} inicializa el periférico I2C
del ESP32 asignando GPIO~21 como SDA y GPIO~22 como SCL. Esta función
configura internamente el driver I2C del ESP32, habilita los pull-ups
internos y establece la velocidad de reloj a 100~kHz por defecto.

\subsubsection{Inicio de Transmisión}
\texttt{Wire.beginTransmission(0x68)} coloca en la cola interna de
transmisión la dirección del esclavo (0x68) con el bit de escritura.
La trama física no se envía hasta llamar a \texttt{Wire.endTransmission()}.

\subsubsection{Escritura de la Dirección del Registro}
\texttt{Wire.write(0x75)} añade a la cola el byte de dirección del
registro \texttt{WHO\_AM\_I}. Este byte le indica al MPU6050 desde qué
posición del mapa de registros se comenzará a leer.

\subsubsection{REPEATED START}
\texttt{Wire.endTransmission(false)} envía la trama de escritura
acumulada y genera una condición de REPEATED START (reinicio sin
liberar el bus), que es la secuencia estándar para cambiar de modo
escritura a modo lectura en I2C sin ceder el control del bus a otro
maestro.

\subsubsection{Solicitud de Lectura}
\texttt{Wire.requestFrom(0x68, 1)} genera la trama de lectura con la
dirección del esclavo y el bit R=1, y solicita la recepción de 1~byte.
El sensor responde con el valor almacenado en el registro \texttt{0x75}.

\subsubsection{Lectura y Presentación del Resultado}
\texttt{Wire.available()} verifica que haya al menos un byte en el
buffer de recepción. Si es así, \texttt{Wire.read()} extrae el byte y
se imprime en el Monitor Serie en formato hexadecimal mediante
\texttt{Serial.println(val, HEX)}.

\subsection{Procedimiento de Ejecución}

\begin{enumerate}
  \item Realizar las conexiones según la Figura~\ref{fig:conexion} y
        la Tabla~\ref{tab:pins}.
  \item Abrir el archivo \texttt{.ino} en Arduino IDE.
  \item En \textit{Herramientas $\rightarrow$ Placa}, seleccionar
        \textbf{ESP32 Dev Module}.
  \item Seleccionar el puerto COM correspondiente al ESP32.
  \item Compilar y cargar (\textit{Ctrl+U}).
  \item Abrir el Monitor Serie (\textit{Ctrl+Mayús+M}) configurado a
        \textbf{115\,200~baudios}.
  \item Observar el valor de \texttt{WHO\_AM\_I} reportado.
\end{enumerate}

% -----------------------------------------------------------------------
%  5. RESULTADOS E IMPLEMENTACIÓN
% -----------------------------------------------------------------------
\section{Resultados e Implementación}
\label{sec:resultados}

\subsection{Resultado Experimental}

Al ejecutar la herramienta sobre el hardware del Nodo~2 del proyecto
(ESP32 DevKit V1 con módulo MPU6050 del mercado local colombiano), la
salida obtenida en el Monitor Serie fue:

\begin{lstlisting}[style=output,
  caption={Salida del Monitor Serie (115\,200 baudios)},
  label=lst:out]
Leyendo WHO_AM_I...
WHO_AM_I: 0x72
\end{lstlisting}

Este resultado confirma simultáneamente tres aspectos del sistema:
\begin{enumerate}
  \item El bus I2C opera correctamente en los pines GPIO~21/22.
  \item El sensor MPU6050 responde a la dirección \texttt{0x68}.
  \item La revisión de silicio del módulo es \texttt{0x72} (no el
        estándar \texttt{0x68}).
\end{enumerate}

\subsection{Tabla de Diagnóstico y Resultados Posibles}

La Tabla~\ref{tab:resultados} clasifica todos los escenarios de salida
posibles y la acción recomendada en cada caso.

\begin{table}[H]
\centering
\caption{Tabla de diagnóstico según valor de WHO\_AM\_I}
\label{tab:resultados}
\scriptsize
\begin{tabular}{p{0.9cm}p{1.2cm}p{3.6cm}}
\toprule
\textbf{Valor} & \textbf{Estado} & \textbf{Diagnóstico y acción} \\
\midrule
0x72 & \textcolor{green!60!black}{\bfseries OK} & Sensor identificado (rev. del proyecto). Ajustar máscara WHO\_AM\_I en driver IDF a 0x72. \\
0x68 & \textcolor{green!60!black}{\bfseries OK} & Sensor original datasheet. Compatible sin modificaciones en el driver. \\
0x70 & \textcolor{orange}{\bfseries ALERTA} & Posible MPU6500. Verificar compatibilidad del driver. \\
0x98 & \textcolor{orange}{\bfseries ALERTA} & Revisión alternativa. Revisar lógica de inicialización. \\
Otro & \textcolor{red}{\bfseries ERROR} & ID desconocido. Confirmar alimentación y cableado I2C. \\
(nada) & \textcolor{red}{\bfseries ERROR} & Sin respuesta. Verificar conexiones, VCC=3.3V y AD0. \\
\bottomrule
\end{tabular}
\end{table}

\subsection{Análisis del Flujo de Comunicación I2C}

La secuencia de tramas I2C generadas por el programa puede
representarse de manera simplificada como:

\begin{lstlisting}[style=output, numbers=none, caption={Trama I2C generada}]
[START][0xD0 W][ACK][0x75][ACK]
[RSTART][0xD1 R][ACK][DATO=0x72][NACK][STOP]
\end{lstlisting}

donde \texttt{0xD0} = dirección 0x68 desplazada 1 bit a la izquierda con bit W=0, y \texttt{0xD1} = dirección 0x68 con bit R=1. El tiempo total de la transacción a 100~kHz es de aproximadamente 180~$\mu$s.

\subsection{Implicaciones para el Firmware ESP-IDF}

La identificación de \texttt{WHO\_AM\_I = 0x72} tiene implicaciones
directas en el firmware de producción del Nodo~2, desarrollado con el
framework ESP-IDF:

\begin{itemize}
  \item La función \texttt{mpu6050\_init()} del driver de bajo nivel
        (archivo \texttt{esp32\_2/main/mpu6050.c} en el repositorio)
        debe modificar su condición de verificación para aceptar el
        valor \texttt{0x72} como identidad válida del sensor.
  \item Las mascaras de bits que filtran el registro devuelto deben
        adaptarse al valor experimental obtenido.
  \item La documentación del README~\cite{readme_debug} registra este
        hallazgo para garantizar la trazabilidad del hardware utilizado
        en experimentos futuros.
\end{itemize}

\subsection{Archivos del Proyecto en GitHub}

La Tabla~\ref{tab:archivos} lista los archivos del repositorio público
relevantes para este artículo.

\begin{table}[H]
\centering
\caption{Archivos del proyecto en el repositorio GitHub}
\label{tab:archivos}
\scriptsize
\begin{tabular}{p{3.5cm}p{3.0cm}}
\toprule
\textbf{Archivo} & \textbf{Descripción} \\
\midrule
\texttt{herramienta\_who\_i\_am\_...ino} & Código fuente principal \\
\texttt{README.txt} & Notas de depuración \\
\texttt{build/.../embebido.ino.bin} & Binario compilado \\
\texttt{build/.../embebido.ino.elf} & ELF con símbolos de depuración \\
\texttt{esp32\_2/main/mpu6050.c} & Driver C (ESP-IDF) \\
\texttt{esp32\_2/main/mpu6050.h} & Cabecera del driver \\
\texttt{esp32\_2/CMakeLists.txt} & Build system IDF \\
\bottomrule
\end{tabular}
\end{table}

% -----------------------------------------------------------------------
%  5b. DISCUSIÓN
% -----------------------------------------------------------------------
\section{Discusión}
\label{sec:discusion}

\subsection{Implicaciones del Valor 0x72}

El hallazgo de \texttt{WHO\_AM\_I = 0x72} en lugar del valor canónico
\texttt{0x68} del datasheet oficial tiene consecuencias técnicas
importantes que van más allá de la simple verificación de identidad.
En primer lugar, confirma que el módulo utilizado es una variante
de terceros —probablemente fabricada por un proveedor chino sin licencia
oficial de TDK— que implementa el mismo núcleo funcional del MPU6050
pero con una programación diferente del registro de identidad. Esta
práctica es conocida en la industria electrónica de componentes genéricos
y no invalida las capacidades funcionales del sensor, pero sí requiere
adaptaciones específicas en el software.

\subsection{Impacto en el Desarrollo del Driver ESP-IDF}

En el contexto del framework ESP-IDF, los drivers de dispositivos
I2C implementan habitualmente una verificación del tipo:

\begin{lstlisting}[style=output, numbers=none]
if (who_am_i != EXPECTED_ID) {
    return ESP_ERR_NOT_FOUND;
}
\end{lstlisting}

Si \texttt{EXPECTED\_ID} está fijado en \texttt{0x68} y el sensor
devuelve \texttt{0x72}, la inicialización falla con el código de error
\texttt{ESP\_ERR\_NOT\_FOUND}, haciendo que el sistema complete el
arranque sin el sensor operativo. La herramienta de diagnóstico
desarrollada permite detectar este problema antes de integrar el driver
completo, ahorrando horas de depuración a nivel de firmware.

\subsection{Generalización a Otros Proyectos}

La metodología propuesta es directamente aplicable a cualquier proyecto
que utilice el MPU6050 u otros sensores I2C en plataformas ESP32 o
similares. La lectura del registro \texttt{WHO\_AM\_I} debería ser un
paso estándar en el flujo de calificación de hardware de cualquier
diseño electrónico que dependa de sensores I2C, especialmente cuando
se adquieren componentes en el mercado secundario o de importación.

\subsection{Limitaciones de la Herramienta}

La herramienta tiene las siguientes limitaciones:
\begin{enumerate}
  \item Solo realiza una lectura al encendido; no monitoriza cambios
        durante la operación.
  \item No verifica el correcto funcionamiento de la medición
        (aceleración/giroscopía), solo la identidad del silicio.
  \item Requiere el entorno Arduino IDE, no es directamente ejecutable
        desde el framework ESP-IDF de producción.
  \item No implementa manejo de errores para bus I2C bloqueado
        (\textit{bus stuck}).
\end{enumerate}

Estas limitaciones son intencionadas en el diseño: la herramienta
es deliberadamente simple para maximizar la fiabilidad del diagnóstico
inicial y minimizar las posibles fuentes de error propias del código.

% -----------------------------------------------------------------------
%  6. CONCLUSIONES Y TRABAJO FUTURO
% -----------------------------------------------------------------------
\section{Conclusiones}
\label{sec:conclusiones}

\subsection{Conclusiones}

\begin{enumerate}
  \item Se diseñó e implementó exitosamente una herramienta de
        diagnóstico de hardware para el sensor MPU6050 sobre ESP32
        DevKit V1 mediante la lectura del registro \texttt{WHO\_AM\_I}
        (dirección \texttt{0x75}) vía I2C, utilizando únicamente la
        biblioteca nativa \texttt{Wire.h} del entorno Arduino IDE, sin
        dependencias externas.

  \item El experimento realizado sobre el hardware del Nodo~2 del
        sistema wearable en la Universidad del Cauca confirmó que el
        módulo MPU6050 utilizado reporta \texttt{WHO\_AM\_I = 0x72},
        que difiere del valor estándar \texttt{0x68} del datasheet
        oficial. Este hallazgo es crítico e imprescindible para la
        correcta configuración del driver de bajo nivel en ESP-IDF.

  \item El diseño minimalista del programa —lectura en \texttt{setup()},
        \texttt{loop()} vacío— garantiza un diagnóstico rápido y sin
        ambigüedades, adecuado como \textbf{Fase~0} obligatoria en
        cualquier flujo de puesta en marcha de hardware basado en
        MPU6050 sobre ESP32.

  \item La configuración hardware documentada (GPIO~21 SDA, GPIO~22 SCL,
        dirección I2C \texttt{0x68}, 115\,200~baudios) fue validada
        experimentalmente y publicada en el repositorio público GitHub,
        facilitando la replicación del entorno en otros nodos del
        sistema y en proyectos similares de la comunidad.

  \item La publicación del código fuente, binarios de compilación
        y documentación de depuración en un repositorio de acceso abierto
        promueve la reproducibilidad de los experimentos y la
        colaboración en el desarrollo de sistemas embebidos wearables
        para monitoreo de AVD.
\end{enumerate}

\subsection{Trabajo Futuro}

Como extensión natural de este trabajo se plantea:
\begin{itemize}
  \item Integrar la verificación de \texttt{WHO\_AM\_I} como
        condición de guarda (\textit{guard clause}) en la función
        \texttt{mpu6050\_init()} del driver ESP-IDF del proyecto.
  \item Ampliar la herramienta para leer y verificar otros registros
        de configuración críticos: \texttt{PWR\_MGMT\_1} (\texttt{0x6B}),
        \texttt{SMPLRT\_DIV} (\texttt{0x19}), \texttt{CONFIG}
        (\texttt{0x1A}) y \texttt{GYRO\_CONFIG} (\texttt{0x1B}).
  \item Desarrollar una versión en ESP-IDF puro (C/FreeRTOS) de la
        herramienta de diagnóstico, coherente con el entorno de
        producción del proyecto wearable.
  \item Diseñar una rutina de autocalibración del offset de
        acelerómetro y giroscopio basada en los datos del registro
        leído, para mejorar la precisión del sistema AVD.
\end{itemize}

% -----------------------------------------------------------------------
%  REFERENCIAS
% -----------------------------------------------------------------------
\begin{thebibliography}{8}

\bibitem{invensense2013}
InvenSense Inc., \textit{MPU-6000 and MPU-6050 Product Specification
Revision 3.4}, doc. PS-MPU-6000A-00, InvenSense Inc., Sunnyvale,
CA, USA, 2013. [En línea]. Disp.:
\url{https://invensense.tdk.com/wp-content/uploads/2015/02/MPU-6000-Datasheet1.pdf}

\bibitem{nxp_i2c}
NXP Semiconductors, \textit{I2C-bus specification and user manual},
Rev.~7.0, UM10204, NXP Semiconductors, 2021. [En línea]. Disp.:
\url{https://www.nxp.com/docs/en/user-guide/UM10204.pdf}

\bibitem{espressif2023}
Espressif Systems, \textit{ESP32 Technical Reference Manual}, v5.1,
Espressif Systems, Shanghai, China, 2023. [En línea]. Disp.:
\url{https://www.espressif.com/sites/default/files/documentation/esp32_technical_reference_manual_en.pdf}

\bibitem{chen2012}
L. Chen, J. Hoey, C. D. Nugent, D. J. Cook y Z. Yu,
``Sensor-based activity recognition,'' \textit{IEEE Trans. Syst.,
Man, Cybern. Part C}, vol.~42, no.~6, pp.~790--808, Nov.~2012,
doi: 10.1109/TSMCC.2012.2198883.

\bibitem{anguita2012}
D. Anguita, A. Ghio, L. Oneto, X. Parra y J. L. Reyes-Ortiz,
``Human activity recognition on smartphones using a multiclass
hardware-friendly support vector machine,'' in \textit{Proc. Int.
Workshop Ambient Assisted Living (IWAAL)}, Vitoria-Gasteiz, Spain,
2012, pp.~216--223.

\bibitem{readme_debug}
C. D. Cuellar Antury, ``Notas de Depuración de Hardware — Proyecto
AVD,'' \textit{README.txt}, Univ. del Cauca, 2024. [En línea]. Disp.:
\url{https://github.com/Cdcuellar/Sistemas-embebidos---Actividades-de-la-vida-diaria/blob/main/herramienta_who_i_am_id_interno_arduino_ide/README.txt}

\bibitem{repo_github}
C. D. Cuellar Antury, ``Sistemas Embebidos — Actividades de la Vida
Diaria,'' Repositorio GitHub, Univ. del Cauca, 2024. [En línea].
Disp.: \url{https://github.com/Cdcuellar/Sistemas-embebidos---Actividades-de-la-vida-diaria}

\bibitem{wire_arduino}
Arduino LLC, \textit{Wire Library Reference}, Arduino Reference
Documentation, 2023. [En línea]. Disp.:
\url{https://www.arduino.cc/en/reference/wire}

\bibitem{kolban2017}
N. Kolban, \textit{Kolban's Book on ESP32}, Leanpub, 2017. [En línea].
Disp.: \url{https://leanpub.com/kolban-ESP32}

\end{thebibliography}

\end{document}
